% Part: first-order-logic
% Chapter: syntax-and-semantics
% Section: substitution

\documentclass[../../../include/open-logic-section]{subfiles}

\begin{document}

\olfileid{fol}{syn}{sub}

\olsection{પ્રતિસ્થાપન}

\begin{defn}[પદમાં પ્રતિસ્થાપન]
$s$માંની $x$ની દરેક ઘટનાના સ્થાને $t$ \emph{પ્રતિસ્થાપિત કરવાથી}
મળતું પરિણામ $\Subst{s}{t}{x}$ નીચે મુજબ પુનરાવર્તી રીતે વ્યાખ્યાયિત
કરીએ છીએ:
\begin{enumerate}
\item \indcase{s}{c}{$\Subst{\indfrm}{t}{x}$ માત્ર $s$ જ છે.}

\item \indcase{s}{y}{$\Subst{\indfrm}{t}{x}$ પણ માત્ર~$s$ જ છે,
  જો $y$ ચલ હોય અને $y \not\ident x$ હોય.}

\item \indcase{s}{x}{$\Subst{\indfrm}{t}{x}$ એ~$t$ છે.}

\item\indcase{s}{\Atom{f}{t_1, \dots, t_n}}{$\Subst{\indfrmp}{t}{x}$ એ
  $\Atom{f}{\Subst{t_1}{t}{x}, \dots, \Subst{t_n}{t}{x}}$ છે.}
\end{enumerate}
\end{defn}

\begin{defn}
પદ~$t$, $!A$માં $x$ માટે \emph{!!{free for}} ત્યારે છે જ્યારે
$!A$માંની $x$ની કોઈ પણ મુક્ત ઘટના એવા પરિમાણકના વ્યાપમાં ન આવે જે
$t$માંના કોઈ ચલને બાંધે છે.
\end{defn}

\begin{ex} ~
\begin{enumerate}
\item $\Obj v_8$, $\lexists[\Obj
  v_3]\Atom{\Obj A^2_4}{\Obj v_3,\Obj v_1}$માં $\Obj v_1$ માટે
  મુક્તપણે પ્રતિસ્થાપનીય છે.

\item $\Obj f^2_1(\Obj v_1, \Obj v_2)$,
  $\lforall[\Obj v_2]\Atom{\Obj A^2_4}{\Obj v_0,\Obj v_2}$માં
  $\Obj v_0$ માટે \emph{મુક્તપણે પ્રતિસ્થાપનીય નથી}.
\end{enumerate}
\end{ex}

\begin{defn}[!!a{formula}માં પ્રતિસ્થાપન]
જો $!A$ !!a{formula} હોય, $x$~!!a{variable} હોય અને $t$, $!A$માં
$x$ માટે !!{free for} પદ હોય, તો $\Subst{!A}{t}{x}$ એ $!A$માંની
$x$ની બધી મુક્ત ઘટનાઓના સ્થાને $t$ પ્રતિસ્થાપિત કરવાથી મળતું
પરિણામ છે.
\begin{enumerate}
\tagitem{prvFalse}{\indcase{!A}{\lfalse}{$\Subst{\indfrm}{t}{x}$ એ
    $\lfalse$ છે.}}{}

\tagitem{prvTrue}{\indcase{!A}{\ltrue}{$\Subst{\indfrm}{t}{x}$ એ
    $\ltrue$ છે.}}{}

\item \indcase{!A}{\Atom{P}{t_1,\dots,
    t_n}}{$\Subst{\indfrm}{t}{x}$ એ $\Atom{P}{\Subst{t_1}{t}{x},
    \dots, \Subst{t_n}{t}{x}}$ છે.}

\item \indcase{!A}{\eq[t_1][t_2]}{$\Subst{\indfrmp}{t}{x}$ એ
  $\Subst{t_1}{t}{x} = \Subst{t_2}{t}{x}$ છે.}

\tagitem{prvNot}{\indcase{!A}{\lnot !B}{$\Subst{\indfrmp}{t}{x}$ એ
    $\lnot \Subst{!B}{t}{x}$ છે.}}{}

\tagitem{prvAnd}{\indcase{!A}{(!B \land
    !C)}{$\Subst{\indfrmp}{t}{x}$ એ $(\Subst{!B}{t}{x} \land
    \Subst{!C}{t}{x})$ છે.}}{}

\tagitem{prvOr}{\indcase{!A}{(!B \lor
    !C)}{$\Subst{\indfrmp}{t}{x}$ એ $(\Subst{!B}{t}{x} \lor
    \Subst{!C}{t}{x})$ છે.}}{}

\tagitem{prvIf}{\indcase{!A}{(!B \lif
    !C)}{$\Subst{\indfrmp}{t}{x}$ એ $(\Subst{!B}{t}{x} \lif
    \Subst{!C}{t}{x})$ છે.}}{}

\tagitem{prvIff}{\indcase{!A}{(!B \liff
    !C)}{$\Subst{\indfrmp}{t}{x}$ એ $(\Subst{!B}{t}{x} \liff
    \Subst{!C}{t}{x})$ છે.}}{}

\tagitem{prvAll}{
\indcase{!A}{\lforall[y][!B]}{$\Subst{\indfrmp}{t}{x}$ એ
    $\lforall[y][\Subst{!B}{t}{x}]$ છે, જો $y$ ચલ હોય અને $x$થી
    જુદો હોય; નહિ તો $\Subst{\indfrmp}{t}{x}$ માત્ર $\indfrm$ જ છે.}}{}

\tagitem{prvEx}{\indcase{!A}{\lexists[y][!B]}{$\Subst{\indfrmp}{t}{x}$ એ
    $\lexists[y][\Subst{!B}{t}{x}]$ છે, જો $y$ ચલ હોય અને $x$થી
    જુદો હોય; નહિ તો $\Subst{\indfrmp}{t}{x}$ માત્ર $\indfrm$ જ છે.}}{}
\end{enumerate}
\end{defn}

\begin{explain}
નોંધો કે પ્રતિસ્થાપન નિરર્થક પણ હોઈ શકે છે: જો $x$, $!A$માં
જરાય ન આવતો હોય, તો $\Subst{!A}{t}{x}$ માત્ર~$!A$ જ છે.

$t$, $!A$માં $x$ માટે !!{free for} હોવું જોઈએ એ પ્રતિબંધ નીચેના
જેવા કિસ્સાઓને બાકાત રાખવા જરૂરી છે. જો
$!A \ident \lexists[y][x < y]$ અને $t \ident y$ હોય, તો
$\Subst{!A}{t}{x}$નું પરિણામ $\lexists[y][y < y]$ થાય. આ કિસ્સામાં
મુક્ત ચલ~$y$ પ્રતિસ્થાપન વખતે પરિમાણક~$\lexists[y]$ દ્વારા
``પકડાઈ'' જાય છે, જે ઇચ્છનીય નથી. ઉદાહરણ તરીકે, જ્યારે પણ
$\lforall[x][!B]$ સાચું હોય ત્યારે $\Subst{!B}{t}{x}$ પણ સાચું હોય
એવું આપણે ઇચ્છીએ. પરંતુ $\lforall[x][\lexists[y][x < y]]$ લો
(અહીં $!B$ એ $\lexists[y][x < y]$ છે). આ વાક્ય, દાખલા તરીકે,
પ્રાકૃતિક સંખ્યાઓ વિશે સાચું છે: દરેક સંખ્યા~$x$ માટે તેનાથી મોટી
કોઈ સંખ્યા~$y$ હોય છે. જો આપણે $x$ના સ્થાને $y$નું પ્રતિસ્થાપન
મંજૂર કરીએ, તો $\Subst{!B}{y}{x} \ident \lexists[y][y < y]$ મળે,
જે અસત્ય છે. $t$માંનો કોઈ પણ મુક્ત ચલ $!A$માંના કોઈ પરિમાણક દ્વારા
બદ્ધ ન થઈ જાય એવી શરત રાખીને આપણે આ અટકાવીએ છીએ.
\end{explain}

અટપટી સંકેતપદ્ધતિ ટાળવા આપણે ઘણી વાર નીચેનો રિવાજ વાપરીએ છીએ: જો
$!A$ એવું !!a{formula} હોય જેમાં !!{variable}~$x$ મુક્ત આવી શકે, તો
આ દર્શાવવા આપણે $!A(x)$ પણ લખીએ છીએ. કયા $!A$ અને~$x$ની વાત છે તે
સ્પષ્ટ હોય અને $t$ પદ હોય ($!A(x)$માં $x$ માટે મુક્તપણે
પ્રતિસ્થાપનીય હોવાનું ધારેલું), ત્યારે $\Subst{!A}{t}{x}$ માટે ટૂંકમાં
$!A(t)$ લખીએ છીએ. તેથી, દાખલા તરીકે, આપણે ``$!A(t)$ને
$\lforall[x][!A(x)]$નું દૃષ્ટાંત
કહીએ છીએ'' એમ કહી શકીએ. આનો અર્થ એ છે કે જો $!A$ કોઈ !!{formula},
$x$ કોઈ !!a{variable}, અને $t$, $!A$માં $x$ માટે મુક્તપણે
પ્રતિસ્થાપનીય પદ હોય, તો $\Subst{!A}{t}{x}$ એ
$\lforall[x][!A]$નું દૃષ્ટાંત છે.

\end{document}
